\chapter{Plague}
\index{Plague}

\section*{Definition and Overview}
Plague is a severe, zoonotic infectious disease caused by the Gram-negative, facultative anaerobic coccobacillus bacterium \textit{Yersinia pestis}. Primarily a disease of wild rodents, it is transmitted to humans and other mammals via the bites of infected rodent fleas (most notably the oriental rat flea, \textit{Xenopsylla cheopis}) or through direct contact with infected animal tissues. Plague manifests in three main clinical forms depending on the route of exposure: bubonic, septicaemic, and pneumonic. Bubonic plague is the most common form, characterized by painful, swollen lymph nodes known as buboes. Pneumonic plague is the most severe and highly contagious form, capable of person-to-person transmission via aerosolised droplets. It is uniformly fatal if untreated, representing a significant public health threat and a class A biosecurity agent.

\section*{Epidemiology}
Plague has a historically devastating record, responsible for three major global pandemics, including the "Black Death" in the 14th century, which decimated a third of the European population. Today, plague remains endemic in several regions of the world, with the highest burden of disease in sub-Saharan Africa, particularly Madagascar, the Democratic Republic of the Congo, and Peru. Sporadic cases also occur in the western United States and parts of Central Asia. The disease is absent from some countries, where strict quarantine regulations are maintained. Globally, between 1,000 and 3,000 cases are reported annually. Outbreaks are often linked to environmental changes, rodent population surges, and poor sanitation in rural or semi-urban areas.

\section*{Aetiology and Pathophysiology}
The pathogenesis of \textit{Yersinia pestis} involves complex vector-host interactions and highly effective immune evasion mechanisms.
\begin{itemize}
    \item \textbf{Vector transmission:} Fleas ingest blood from an infected rodent. The bacteria multiply inside the flea's proventriculus, forming a biofilm block. This forces the "blocked" flea to bite hosts repeatedly, regurgitating bacteria into the bite wound.
    \item \textbf{Intracellular survival:} Once inside the mammalian host, \textit{Y. pestis} is engulfed by macrophages, where it survives and replicates inside the phagosome, developing a protective capsule that prevents subsequent phagocytosis.
    \item \textbf{Type III secretion system:} The bacterium utilizes a needle-like type III secretion system to inject outer proteins (Yops) directly into host immune cells, disrupting the cytoskeleton, preventing phagocytosis, and inhibiting pro-inflammatory cytokine release.
    \item \textbf{Lymphatic dissemination (Bubonic):} Bacteria travel via the lymphatics to regional lymph nodes, causing severe hemorrhagic necrosis and extreme swelling (bubo formation).
    \item \textbf{Haematogenous spread (Septicaemic/Pneumonic):} Bacteria enter the bloodstream, causing systemic inflammatory response syndrome (SIRS), disseminated intravascular coagulation (DIC), tissue ischemia, and acral gangrene (necrosis of fingers and toes).
\end{itemize}

\section*{Clinical Features}
The incubation period ranges from 1 to 7 days, but can be as short as 24 hours for the pneumonic form. Symptoms vary by clinical subtype:
\begin{itemize}
    \item \textbf{Bubonic plague:} Sudden onset of high fever, chills, headache, and severe malaise. This is followed by the rapid development of extremely painful, erythematous, swollen lymph nodes (buboes) in the groin, armpit, or neck.
    \item \textbf{Septicaemic plague:} Presents with fever, rigors, abdominal pain, and rapid progression to septic shock, hypotension, and organ failure. Haemorrhagic manifestations include petechiae, ecchymoses, and gangrene of the extremities. Buboes may be absent.
    \item \textbf{Pneumonic plague:} Rapidly progressive necrotising pneumonia presenting with high fever, severe cough, dyspnoea, tachypnoea, and bloody, frothy sputum.
    \item \textbf{Neurological signs:} Agitation, confusion, meningitis, or coma can occur in late-stage systemic infection.
\end{itemize}

\section*{Differential Diagnosis}
Given its non-specific initial symptoms, plague must be differentiated from other acute infectious diseases.
\begin{itemize}
    \item \textbf{Severe tularaemia:} Caused by \textit{Francisella tularensis}; presents with ulceroglandular features and lymphadenopathy, but typically lacks the extreme pain and rapid progression of plague buboes.
    \item \textbf{Bacterial lymphadenitis:} Staphylococcal or streptococcal infections causing localized lymph node swelling; distinguished by purulence and microbiological culture.
    \item \textbf{Severe community-acquired pneumonia:} Pneumonias caused by \textit{Streptococcus pneumoniae} or \textit{Legionella}; distinguished by clinical history (lack of travel to plague-endemic regions) and sputum PCR.
    \item \textbf{Meningococcemia:} Presents with rapid sepsis and purpuric skin lesions; distinguished by blood culture or PCR.
    \item \textbf{Scrub typhus or rickettsial diseases:} Vector-borne diseases presenting with fever and eschars; diagnosed via serology.
\end{itemize}

\section*{When to Seek Medical Care}
Any individual who develops a sudden fever and painful, swollen lymph nodes or severe respiratory symptoms after travelling to a plague-endemic region must seek immediate, isolated medical care.

\textbf{Contact your local emergency services immediately for:}
\begin{itemize}
    \item Severe difficulty breathing, rapid respiration, or coughing up blood
    \item High fever with confusion, extreme lethargy, or loss of consciousness
    \item Sudden, extremely painful swelling of the lymph nodes in the groin, armpit, or neck
    \item Signs of shock, such as cold, clammy skin, rapid heart rate, and low blood pressure
    \item Uncontrolled bleeding or black discolouration of the fingers or toes
\end{itemize}

\section*{Diagnosis and Investigations}
Prompt diagnostic testing is essential, but treatment must not be delayed while waiting for laboratory confirmation.
\begin{itemize}
    \item \textbf{Microbiological culture:} Isolation of \textit{Y. pestis} from bubo aspirates, blood, or sputum. The organism exhibits a characteristic "safety-pin" appearance (bipolar staining) on Wayson or Giemsa stain. The laboratory must be notified of suspected plague to ensure Biosafety Level 3 (BSL-3) containment.
    \item \textbf{Polymerase Chain Reaction (PCR):} Rapid detection of \textit{Y. pestis} DNA in blood, sputum, or tissue samples.
    \item \textbf{Antigen detection:} Enzyme-linked immunosorbent assay (ELISA) or lateral flow strips to detect the F1 capsule antigen.
    \item \textbf{Supportive laboratory findings:} Leukocytosis with a left shift, thrombocytopenia, elevated inflammatory markers (CRP), and signs of DIC or renal failure.
\end{itemize}

\section*{Management}
Plague is highly responsive to antibiotics if initiated early, preferably within 24 hours of symptom onset.
\begin{itemize}
    \item \textbf{Antimicrobial therapy:} First-line agents include aminoglycosides (intravenous streptomycin or gentamicin). Fluoroquinolones (ciprofloxacin or levofloxacin) and tetracyclines (doxycycline) are highly effective alternatives.
    \item \textbf{Isolation precautions:} Patients with suspected pneumonic plague must be isolated in single rooms with negative pressure. Healthcare staff must adhere to droplet precautions, including using fit-tested N95 respirators, until the patient has completed 48 hours of effective antibiotic therapy.
    \item \textbf{Post-exposure prophylaxis (PEP):} Close contacts of patients with pneumonic plague should receive oral doxycycline or ciprofloxacin for 7 days to prevent infection.
    \item \textbf{Supportive care:} Intensive care management for patients in septic shock or respiratory failure, including intravenous fluid resuscitation, vasopressors, and mechanical ventilation.
\end{itemize}

\section*{Complications}
\begin{itemize}
    \item Septic shock and progressive multi-organ failure.
    \item Disseminated intravascular coagulation (DIC) leading to severe haemorrhage and gangrene.
    \item Adult respiratory distress syndrome (ARDS) secondary to pneumonic plague.
    \item Plague meningitis, a rare complication due to haematogenous seeding of the meninges.
    \item High rate of miscarriage and maternal mortality in pregnant patients.
\end{itemize}

\section*{Prognosis and Outcomes}
The prognosis of plague depends entirely on the timing of antibiotic initiation. With early, appropriate antibiotic treatment, the mortality rate for bubonic plague is reduced from approximately 50\%--60\% to less than 10\%. In contrast, septicaemic and pneumonic plagues carry a mortality rate of virtually 100\% if untreated, and if treatment is delayed by more than 24 hours after symptom onset, the mortality rate remains high. Survivors typically recover fully without long-term sequelae, although tissue loss may occur in patients who develop severe acral gangrene.

\section*{Prevention and Self-Care}
\begin{itemize}
    \item Reduce rodent habitats around homes by clearing trash, woodpiles, and food waste.
    \item Protect pets from fleas using veterinary-approved flea control products, and keep them indoors in endemic areas.
    \item Use insect repellent containing DEET or picaridin when camping, hiking, or working outdoors in plague-endemic regions.
    \item Avoid direct contact with dead animals or rodents in endemic zones.
    \item Adhere to recommended post-exposure prophylactic antibiotic courses if exposed to a confirmed case.
\end{itemize}

\section*{Sources and Further Reading}
The following reputable organisations provide detailed clinical guidelines and public health resources on plague:
\begin{itemize}
    \item World Health Organization. \textit{Plague fact sheets, epidemiology, and global surveillance.} https://www.who.int/
    \item Centers for Disease Control and Prevention. \textit{Plague clinical overview, diagnosis, and treatment protocols.} https://www.cdc.gov/
    \item European Centre for Disease Prevention and Control. \textit{Factsheet on plague for health professionals.} https://www.ecdc.europa.eu/
    \item Australian Government Department of Health and Aged Care. \textit{Biosecurity and quarantine protocols for high-consequence pathogens.} https://www.health.gov.au/
    \item Healthdirect Australia. \textit{Infectious diseases and international travel health advice.} https://www.healthdirect.gov.au/
\end{itemize}

\clearpage
